\documentclass[12pt]{article}
\usepackage[utf8]{inputenc}
\usepackage[T1]{fontenc}
\usepackage{hyperref}
\title{Southeastern Tarahumara Language Resources for K-12 Curriculum Development}
\author{Compiled by Claude Sonnet 5 (claude-sonnet-5) for Andr\'as Kornai}
\date{}
\begin{document}
\maketitle

\section{Introduction}
Southeastern Tarahumara (ISO 639-3 tcu; Glottocode sout2972) is an endangered Uto-Aztecan language of Mexico. The reference description says it is used as a first language by adults only and is not known to be taught in schools. Web search found no resource specific to this variety. Most of what exists covers Tarahumara (Rar\'amuri) in general or other varieties. Check the dialect of any item before using it for tcu.

\section{Primary Educational Resources}
No primers, textbooks or school materials for Southeastern Tarahumara turned up in search.

\section{Linguistic and Reference Materials}
\begin{itemize}
\item Tarahumarisches W\"orterbuch (1791), a German--Tarahumara and Tarahumara--German dictionary held by the University of Arizona Special Collections. It is historical and not specific to the southeastern variety: \url{https://lib.arizona.edu/special-collections/collections/tarahumarisches-worterbuch-nebst-einigen-nachrichten-von-den-sitten}
\item A Princeton catalogue record for a grammatical description of Choguita Rar\'amuri. This is a different variety, and the search snippet says it is backed by over 200 hours of documentary recordings in two archives: \url{https://catalog.princeton.edu/catalog/99126844766406421}
\item Tarahumara materials in the Jane M. Rosenthal Papers (American Philosophical Society), an archival finding aid: \url{https://indigenousguide.amphilsoc.org/entry/39272}
\item A 1955 Tarahumara--English and English--Tarahumara dictionary by I. Thord-Gray appears in a bookseller listing. It covers Tarahumara generally: \url{https://www.abaa.org/book/867166792}
\item Wikipedia's page for ISO 639:tcu, useful for the code and variety classification: \url{https://wikipedia.com/wiki/ISO_639:tcu}
\end{itemize}

\section{Religious and Community Resources}
None found in search.

\section{Digital and Audio Resources}
\begin{itemize}
\item Endangered Languages Project records for Tarahumara varieties. I did not verify which variety each record covers: \url{https://endangeredlanguages.com/node/96676} and \url{https://endangeredlanguages.com/node/38913}
\end{itemize}

\section{Curriculum Development Context}
Little exists for tcu specifically, and no school programme was found. A curriculum would have to start from new documentation. Contacting Mexican institutions such as INALI and local communities is a sensible next step, but this search did not surface specific contacts.
\end{document}